\section{Feral Spren Encounters}
\label{ssec:shadesmarFeralEncounters}

Strong emotions attract spren in the Cognitive Realm just as as usual. However, in here the spren can physically interact with the characters and even attack them.

When a PC shows an emotion, either during a player's roleplay or an action taken by the character, you can consider an encounter similar to one of the suggested here: 



\subsection{Fearspren}

Scared PCs can attract \enemyFearSpren{}s. They converge towards the character subjected to fear, and unless the PC succeeds on \skillCheck{14}{\skillDiscipline{}} test to calm down, they will engage.

If the character fails the test, 2d4 \enemyFearSpren{} will attack them. They are not lethal, but they can harm them in other ways. During the fight, the character can repeat the test with disadvantage as~\iconAction{}.

The fight ends either when the spren are defeated, or when the character successfully passes the test. The spren will also retreat after they successfully attack a character 4 times.

\spoilerBox{
	Adolin attracting fearspren is shown in \bookOB{}, Chapter 89. Damnation.
}



\subsection{Angerspren}

A character filled with anger can attract \enemyAngerSpren{}s. At first, the party can hear only screeching, metallic
noises of the pack. When they do, allow the angry PC for a \skillCheck{14}{\skillDiscipline{}} while they attempt to calm down.

If they fail, they are attacked by 1d4 \enemyAngerSpren{}s. They can repeat the check during the fight as~\iconAction{}, but now it is made at disadvantage.

The fight ends either when angerspren are defeated or when the angry character manages to calm down.

\spoilerBox{
	Kaladin's anger attracting the spren is shown in \bookOB{}, Chapter 95. Inescapable Void.
}



\subsection{Painspren}

Characters who are hurt may attract up to 2d6 \enemyPainSpren{}s. They will latch to the hurting PC, causing even more pain.

Those cannot be escaped from, the encounter ends when they are defeated in a fight or when the hurting character looses consciousness. 




\subsection{Anticipationspren}

When a character is nervous or anticipates something to happen, they can attract 1d6 \enemyAnticipationSpren{}s. They are not hostile, but they signal the PC's location with their long, streamer-like tongues, granting disadvantage to any \skillStealth{} checks.

If a character decides to attack one of them, they will defend themselves until defeated.



\subsection{Lifespren}

Lifespren can be found in the parts of the Cognitive Realm that correspond to the lush areas of the Physical Realm, overgrown with verdant life.

Those spren are not interested with the PCs and they will continue to swarm around the plants they detect in the Physical Realm. Should a player gently interact with them, however, they will heal them for 1d4 Health. This can happen only once per encounter.

\reasonBox{
	Lifespren were never shown on a cards of \bookSA{} in a scene where they healed someone, but I think it's good to have something positive happening to the party in such a hostile environment.
}



\subsection{Creationspren}

Creationspren can be attracted by a PC using their \skillCrafting{} skill, eg. during a free time. They are attracted to the process of creation and they are intrigued by the thing being crafted.

At the end of this encounter, the character attracting them receives 1 Focus and 1 Maximum Focus until the next Long Rest.

